\needspace{10\baselineskip}
\hypertarget{native-322-zusatz17}{}
\pdfbookmark[1]{17. Zodiac on the celestial sphere}{zusatz-017}
\begin{center}
[17.]

Zodiac on the celestial sphere.

Oct. 28, 1847.
\end{center}

For every point on the celestial sphere, there may be no combination, one combination, or more than one combination that makes one planet visible at that point from the other, that is, there is a definite number of solutions. Each time a limit line is crossed, the number of solutions changes by 2. As long as the limit line is not encountered along any given path on the celestial sphere, however, the number of solutions remains unchanged. The question now is how to distinguish the side of the limit with two more solutions from the other side with two fewer. The following considerations serve this purpose.

1. Let \(p\) be a point of the planetary orbit from which the points of the other orbit are projected onto the celestial sphere, \(P\) the \SourceLeaf{323}corresponding point of the other orbit, and \(\Pi\) the point on the line of nodes where the tangents to the two planetary orbits at \(p\) and \(P\) intersect. Further, let
\[
p\Pi=n,\qquad P\Pi=N,\qquad pP=\Delta,
\]
\(S\) be any other fixed point on the line of nodes (for example, the Sun). The three figures are to be imagined in three different planes: I in the plane through \(\Pi pP\), II in the plane through \(\Pi Sp\), and III in the plane through \(\Pi SP\). We further denote by
\begin{quote}
\noindent \(p,P\) the angles of planes II and III with I,\\
\(q,Q\) the angles \(p\Pi S,P\Pi S\),\\
\(u,U\) the angles \(\Pi pP,\Pi Pp\),\\
\(r,R\) the radii of curvature of the two orbits at \(p\) and \(P\).
\end{quote}

Now let \(p^{\prime},P^{\prime}\) be two other corresponding points, with \(\Pi^{\prime}\) the associated point on the line of nodes. Assume \(\Pi\Pi^{\prime}=\delta\) to be infinitesimal. Then the perpendiculars from \(p^{\prime},P^{\prime}\) to \(pP\) are:
\[
p^{*}p^{\prime}=\frac{\delta r\sin q\sin u}{n}\mathpunct{\raisebox{.45ex}{\normalfont,}}\qquad
P^{*}P^{\prime}=\frac{\delta R\sin Q\sin U}{N}\mathpunct{\cdot}
\]

Further, let \(p^{\prime\prime},P^{\prime\prime}\) be two other points on the tangents \(p\Pi,P\Pi\), such that \(p^{\prime\prime}P^{\prime\prime}\) is parallel to \(pP\) and \(\frac{pp^{\prime\prime}}{p\Pi}=\frac{PP^{\prime\prime}}{P\Pi}=\omega\) is infinitesimal. Let a plane through \(p^{\prime\prime}P^{\prime\prime}\) perpendicular to I meet the curves at \(p^{\prime\prime\prime},P^{\prime\prime\prime}\), and let \(p^{\mathrm{IV}},P^{\mathrm{IV}}\) be the feet of the perpendiculars from \(p^{\prime\prime\prime},P^{\prime\prime\prime}\) to I. We then have
\[
\begin{aligned}
p^{\prime\prime\prime}p^{\mathrm{IV}}&=p^{\prime\prime}p^{\prime\prime\prime}\mathord{.}\sin p=\frac{nn\omega\omega}{2r}\sin p\\[5pt]
P^{\prime\prime\prime}P^{\mathrm{IV}}&=P^{\prime\prime}P^{\prime\prime\prime}\mathord{.}\sin P=\frac{NN\omega\omega}{2R}\sin P\\[5pt]
p^{\prime\prime\prime}p^{\mathrm{IV}}\mathord{.}p^{*}p^{\prime}&=\frac{\delta\omega\omega n\sin p\sin q\sin u}{2}\\[5pt]
P^{\prime\prime\prime}P^{\mathrm{IV}}\mathord{.}P^{*}P^{\prime}&=\frac{\delta\omega\omega N\sin P\sin Q\sin U}{2}\mathpunct{\cdot}
\end{aligned}
\]

But it is easy to see that \(n\sin u=N\sin U\) is the perpendicular from \(\Pi\) to \(pP\). Furthermore, \(\sin p\sin q=\sin P\sin Q={}\) sine of the angle that the line \(\Pi S\) makes with plane I and hence with the perpendicular just mentioned. We easily conclude that \(n\sin p\sin q\sin u=N\sin P\sin Q\sin U\) is simply the shortest distance between the straight lines \(pP,\Pi S\).

\SourceLeaf{324}The final result is therefore
\[
p^{\prime\prime\prime}p^{\mathrm{IV}}\mathord{.}p^{*}p^{\prime}=P^{\prime\prime\prime}P^{\mathrm{IV}}\mathord{.}P^{*}P^{\prime}.
\]
If, therefore, \(p^{\prime\prime\prime}p^{\mathrm{IV}}>P^{\prime\prime\prime}P^{\mathrm{IV}}\), then \(p^{*}p^{\prime}<P^{*}P^{\prime}\), and conversely. The first inequality indicates that a straight line drawn parallel to \(p^{\prime\prime\prime}P^{\prime\prime\prime}\) from \(p\), on the one hand, and \(S\), on the other, lie on opposite sides of I. The second indicates that the projections onto the celestial sphere of the two straight lines \(p^{\prime}P^{\prime}\) and \(pP^{\prime}\) (the latter lying only a 2\textsuperscript{nd}-order infinitesimal away from the limit line) lie on the same side of the projection of \(pP\).

The desired criterion is easily derived from this and can be stated as follows:

[Let] \(B\) [be the] outer planetary orbit as seen on the celestial sphere from the Sun. This great circle divides the celestial sphere into two halves (1) and (2). Let (1) be the half in which the heliocentric position of the inner planet appears. The region with two more solutions then lies with respect to a segment \(LL^{\prime}\) of the limit in the same way as (1) lies with respect to \(B\), taking corresponding senses as forward along \(LL^{\prime}\) and \(B\). This assumes that the planet in the outer orbit is the one being observed. In the opposite case the relation is reversed.

In symbols the criterion can be expressed more generally as follows:

Let \(l,l^{\prime}\) be two points on a limit infinitely close to one another, \(p,p^{\prime}\) the corresponding heliocentric positions of the planet from which [the observation is made], \(P,P^{\prime}\) those of the planet being observed, and \(a\) a point near \(ll^{\prime}\) on the side with two more solutions than the other. Then, if the congruence sign \(\equiv\) [denotes] the same orientation and a prefixed \(-\) sign denotes the opposite orientation, we have:
\[
ll^{\prime}a\equiv PP^{\prime}p\equiv-pp^{\prime}P.
\]
\begin{center}\rule{.22\textwidth}{.3pt}\end{center}
